\scanpage{25}
\deleted{$h\widehat c=\widehat c$}
\diagram{\begin{tikzcd}[column sep=small]
\left\{\begin{array}{c}(P,\sigma)\\\widehat S\end{array}\right.
\arrow[r,leftrightarrow] & 0 \arrow[r,dash] & \Omega \arrow[r,dash] & E \arrow[r] & \mathcal O_S \arrow[r] & 0
\end{tikzcd}}
$\unsure{\mathcal P|P}$ muni \uline{d'une connexion} $C_P$
\figure{\begin{tikzpicture}[x=1cm,y=1cm]
\draw[->] (0,1.7) .. controls (-.6,1) and (-.4,.7) .. (-.1,.4) .. controls (.1,0) and (.3,0) .. (.7,0) -- (2.9,0) -- (2.9,-.35);
\end{tikzpicture}}
On cherche les iso \inserted{\unsure{étg|etc.}} $\unsure{\mathcal P|P}\xrightarrow{v}\widehat S$ \inserted{(s'il y en a)}, pour lesquels $v(C_P)$ est de la forme $\widehat c$.

La donnée d'un iso. à l'ordre $2$ donne déjà l'extension
\diagram{\begin{tikzcd}[column sep=small]
0 \arrow[r,dash] & \Omega \arrow[r,dash] & E \arrow[r] & \mathcal O_S \arrow[r,dash] & \unsure{0}
\end{tikzcd}}
i.e. définit $(P,\sigma)$ à isom unique près.

Si on se donne \uline{un} isom. entre les $\widehat S$ can. et $\unsure{\mathcal P|P}$, $C_P$ est \unsure{défini} par un $C\in\unsure{\operatorname{Conn}_{\mathrm{proj}}}(\widehat S)$ \inserted{qui prolonge la conn. à l'ordre $2$}. On cherche donc les $g\in G$ (autom \inserted{de $\widehat S$} qui sont l'identité dans l'ordre $2$), tels que $g.C$ de la forme $\widehat c$.

\deleted{\uline{\unsure{Supposons}}} \inserted{Il existe tjrs} \deleted{\unsure{qu'il existe}} un tel $\deleted{\unsure{g_0}}\inserted{\unsure{g'_0|g_0}}$ \deleted{\unsure{à condition}} \inserted{\uline{dès qu'il}} \unsure{y a une connexion} $c_0$ \unsure{sur} $P$, car on aura $gC=h\widehat c_0$
\margin{$(h\in G_0)$ \unsure{unique} \unsure{donc}}
\[
\deleted{g_0C=\widehat c}\quad (h^{-1}g)C=\widehat c_0.
\]
Quels sont les autres ? Ils sont de la forme $g=hg_0$. $(h\in G)$ et la condition est que $gC$ $(=h(g_0C)=h\widehat c)$ soit de la forme $\widehat c'$, i.e. $h\in G_1(c)$.

L'ens de ces $h$ est en correspondance biunivoque avec $\unsure{G/G_0|G_1/G_0}$ -- on peut dire aussi que pour \unsure{tte extension des}
% archive foliation: 24
